\scanpage{81}
% Archive pencil foliation: 80.
de $\handwritten{\underline}{Q}_p$ \handwritten{)} qui sont donc isomorphes \inserted{\handwritten{``géométriquement'' i.e. dans $\operatorname{Rep}(\underline{H})$}}\handwritten{\begin{tikzpicture}[baseline=0cm,x=1cm,y=1cm]\draw[thin] (0,-.42)--(0,0)..controls (0,.12) and (.08,.12)..(.18,.12)--(3.1,.17)--(3.1,.28);\end{tikzpicture}} (mais sans doute non canoniquement) aux $E_{r,1}$ pour $r\geq1$: ils donnent des modules libres de rang 1 canoniqu\unsure{\deleted{d}e|\deleted{e}d}s sous les $K_{r,1}$ ...\handwritten{]}

9. Pour $k$ quelconque, on trouve un \deleted{ho} $\handwritten{\otimes}$-homomorphisme canonique \deleted{(à l'aide} défini à isomorphisme unique près (on utilise un choix d'une clôture algébrique $\handwritten{\overline}{k}$ de $k$, mais ce choix est inessentiel ..)
\[
\begin{tikzcd}[cramped]\operatorname{Fcriso}(k)\arrow[r]&\operatorname{Rep}(\handwritten{\underline}{H})\end{tikzcd}
\]
ou encore un homomorphisme de gerbes
\[
\begin{tikzcd}[cramped]\handwritten{\underline}{H}\arrow[r]&\handwritten{\underline}{G}(k)\end{tikzcd}
\]
(\unsure{\deleted{a}r|\deleted{r}a|r}appelons que cet homomorphisme se factorise par le noyau,\deleted{de} soit \deleted{\unsure{Gl|G1|G|Gr}} $\handwritten{\underline}{H}(k)$, de \begin{tikzcd}[cramped]\handwritten{\underline}{G}(k)\arrow[r]&\handwritten{\underline{\Pi}}_p(k)\end{tikzcd}, et que p\unsure{\deleted{i}o|\deleted{o}i}ur $k$ parfait on devine que l'homomorphisme obtenu \begin{tikzcd}[cramped]\handwritten{\underline}{H}\arrow[r]&\handwritten{\underline}{H}(k)\end{tikzcd} est une équivalence ...). Si $\handwritten{\lambda\in}\handwritten{\underline}{Q}$, un F\deleted{(}\handwritten{-}isocristal sur $k$ est dit \deleted{homogèn} \handwritten{\uline}{de pente} $\handwritten{\lambda}$ si son image dans $\operatorname{Rep}(\handwritten{\underline}{H})$ est de pente $\handwritten{\lambda}$. On voit alors aussitôt qu'un F-isocris\unsure{ta|\deleted{t}a|\deleted{a}t}l est de pente zéro sss il est ``trivial'' i.e. satisfait aux conditions de 5.2 (ce qui justifie la terminologie qui est introduite, un peu prématurément, à cet endroit), il est effectif sss son image dans \inserted{$\operatorname{Rep}(H)$ l'est,i.e. est à pente $\handwritten{\geq}$ 0.}

Considérons l'homomorphisme de liens associé
\[
\begin{tikzcd}[cramped]H\arrow[r]&G(k)\quad ,\end{tikzcd}
\]
on trou\unsure{\deleted{e}v|\deleted{v}e}e al\unsure{\deleted{r}o|\deleted{o}r|o}\unsure{\deleted{n}r|\deleted{r}n|\deleted{o}r|\deleted{r}o|n|r|o}s que \unsure{\deleted{e}s|\deleted{s}e}i $k$ est parfait, l'homomorphisme précédent est \handwritten{\uline}{central}, en d'autres termes $H$ opère sur le foncteur identique de \begin{tikzcd}[cramped]\operatorname{Fcriso}(k)\arrow[r]&\operatorname{Rep}(\handwritten{\underline}{G}(k))\end{tikzcd}, de sorte que tout F-isocristal \inserted{$(E,F_E)$}\handwritten{\begin{tikzpicture}[baseline=0cm,x=1cm,y=1cm]\draw[thin] (0,-.18)..controls (.04,.15) and (.38,.12)..(.65,.08)..controls (.9,.05) and (1.25,.16)..(1.3,.24)..controls (1.33,.27) and (1.3,.3)..(1.28,.32);\end{tikzpicture}} se décompose en ses \handwritten{``}\handwritten{\uline}{composantes isopentiques}\handwritten{''} $E_{\handwritten{\lambda}}$, $\handwritten{\lambda\in}\handwritten{\underline}{Q}$. $E$ est ``trivial'' resp. effectif sss $E_{\handwritten{\lambda}}=0$ pour $\handwritten{\lambda\ne}$ 0 (resp. $\handwritten{\lambda<0}$).

Notons que le \inserted{F-iso}\handwritten{\begin{tikzpicture}[baseline=0cm,x=1cm,y=1cm]\draw[thin] (0,-.2)..controls (0,.13) and (.38,.14)..(.62,.16)..controls (.89,.18) and (1.02,.33)..(.97,.36);\end{tikzpicture}}cristal de Tate $K(1)$ est de pente\handwritten{-}1. Donc (pour $k$ parfait) tout F\deleted{(}\handwritten{-}isoc\unsure{\deleted{i}r|\deleted{r}i}\unsure{\deleted{r}i|\deleted{i}r}stal se met de façon unique sous la forme
\[
E=\handwritten{\sum_{i\in\mathbb{Z}}}E_i(-i)\quad ,
\]
où $E_i$ est à composantes isopentiques de pente $0\handwritten{\unsure{<|\leq}}\,\inserted{\handwritten{\rlap{\raisebox{.5em}{\begin{tikzpicture}[x=1cm,y=1cm]\draw[gray,thin] (-.15,.1)..controls (0,.38) and (.45,.3)..(.55,-.03);\end{tikzpicture}}}i}}\handwritten{\leq}1$. D'ailleurs,

